\subsection{P-14: dense interpolation preserves coefficient rows across rank-one resets}
\paragraph{Statement P-14.}
Let $F$ be a field, let $n,w\geq1$, and consider an ROABP in any
one variable order
\[
 f(x)=uA_1(x_{\pi(1)})\cdots A_n(x_{\pi(n)})v,
 \qquad A_i\in F[X]^{w\times w},\quad
 u\in F^{1\times w},\quad v\in F^{w\times1}.
 \tag{P14.1}
\]
Assume that every factor has generic rank either $1$ or $w$ over
$F(X)$. No bound is placed on the number of rank-one factors.
Choose distinct $a_1,\ldots,a_n\in F$, their Lagrange polynomials
$\lambda_i$, and arbitrary $c_i\in F^*$. Put $K=w-1$ and
\[
 G_i(t,\mathbf y,\mathbf z)
       =c_it+\sum_{j=1}^{K}z_j\lambda_i(y_j).
 \tag{P14.2}
\]
Then $f\ne0$ implies $f(G)\ne0$. The generator does not depend
on the order $\pi$ or on which factors have rank one.

More precisely, for any fixed $u$ let
$R(x)=uA_1(x_{\pi(1)})\cdots A_n(x_{\pi(n)})$ and let
$V\subseteq F^{1\times w}$ be its coefficient-row span. Then
\[
 \operatorname{span}_F\{\text{coefficient rows of }R(G)\}=V.
 \tag{P14.3}
\]
This is equality of coefficient-row spaces. It does not assert
equality of the full matrix image cones.

If $\deg f\leq D$, where $D\geq0$, the generator gives a hitting
set of size at most
\[
 (D+1)\bigl((n-1)D+1\bigr)^{w-1}
             \binom{D+w-1}{w-1}.
 \tag{P14.4}
\]
It suffices to work over a field with at least
$\max\{n,D+1,(n-1)D+1\}$ elements; an extension may be used.
All characteristics are allowed. For $w=1$, the empty sum in
(P14.2) is the diagonal line $x_i=c_it$.

In particular every nonzero width-two ROABP in one hidden order,
of arbitrary individual degree, is hit by a common center plus
one interpolation pair. For $n\geq2$, the common center alone
does not hit all width-two ROABPs.

There is also a rational-function form of the result. For distinct
abscissas $a_i$, the evaluation map
\[
 x_i=\frac{g(a_i)}{h(a_i)},\qquad
 \deg g,\deg h\leq w-1,
 \tag{P14.5}
\]
hits the same scalar class wherever its denominators are defined.
The corresponding rational coefficient vectors preserve the
coefficient-row span after expansion in any common polynomial
denominator. The rational form is not needed to construct the
polynomial hitting set (P14.4).

\paragraph{1. A bridge with invertible interior.}
We first prove an elementary local lemma. Let
\[
 b(X)=L(X_0)B_1(X_1)\cdots B_q(X_q)Q(X_{q+1})
 \tag{P14.6}
\]
be a nonzero scalar polynomial, with separate variables in the
displayed factors. Either endpoint may instead be constant.
Assume that $L,Q$ are nonzero polynomial row and column vectors,
respectively, and every $B_i$ is a generically invertible
$w\times w$ polynomial matrix. At the center $X_j=c_jt$, over
$F(t)$, every interior constant is invertible and both endpoint
constants are nonzero. Here each $c_j$ is nonzero and $t$ is
transcendental; substitution into a nonzero univariate polynomial
therefore remains nonzero.

Write $X_j=c_jt+Y_j$. If $L_0=L(c_0t)$, choose an index $a$ with
$(L_0)_a\ne0$. Then
\[
 L(c_0t+Y)=L_0\left(I+
       \frac{e_a}{(L_0)_a}\bigl(L(c_0t+Y)-L_0\bigr)\right).
 \tag{P14.7}
\]
The matrix in parentheses has constant term $I$. Likewise, if
$Q_0=Q(c_{q+1}t)$ and $(Q_0)_b\ne0$, then
\[
 Q(c_{q+1}t+Y)=\left(I+
       \bigl(Q(c_{q+1}t+Y)-Q_0\bigr)
                      \frac{e_b^{\mathsf T}}{(Q_0)_b}\right)Q_0.
 \tag{P14.8}
\]
Thus endpoints can be absorbed into two additional univariate
matrix factors with invertible constants, without increasing
the width or repeating any variable inside this bridge.

For completeness, a product with all constant factors invertible
can be normalized so that every constant is $I$. If its factors
are $M_i(Y_i)$, put $C_i=M_i(0)$, $P_i=C_1\cdots C_i$, $P_0=I$,
and
\[
 N_i(Y_i)=P_{i-1}M_i(Y_i)P_i^{-1}.
\]
Then $N_i(0)=I$ and
$N_1\cdots N_s=M_1\cdots M_sP_s^{-1}$; absorb $P_s$ into
the final constant column.

Let $W_i$ be the coefficient-row span of
$L_0N_1\cdots N_i$ over $F(t)$. Because $N_i(0)=I$,
$W_{i-1}\subseteq W_i$. Let $g_i$ count the strict inclusions
up to position $i$. Induction shows that $W_i$ is spanned by
coefficient rows whose monomials use at most $g_i$ of the
variables. If there is no growth, use the already available
spanning rows with exponent zero in $Y_i$. If there is growth,
multiply the preceding spanning rows by coefficient matrices
of $N_i$; this uses at most one additional variable. Since
$\dim W_0=1$ and $\dim W_i\leq w$, one has $g_i\leq w-1$.
Consequently a nonzero scalar contraction of the final row
has a nonzero coefficient supported on at most $w-1$ variables.

The standard interpolation map with $w-1$ pairs hits any
polynomial with such a coefficient. Indeed, set its selectors
to the abscissas of that support and set any unused amplitudes
to zero. The map then keeps the chosen coordinates as independent
amplitudes and sets every other residual coordinate to zero.
The resulting coordinate restriction is nonzero. This proves
that $b(G)\ne0$, allowing the variables absent from the bridge
to retain their original indices among all $n$ Lagrange nodes.

\paragraph{2. Exact splitting at every rank-one factor.}
Return to (P14.1), and first suppose there is at least one rank-one
factor. Denote its positions by $i_1<\cdots<i_s$ in the program
order. For each such factor choose a nonzero entry $a_j(X)$,
in row $p_j$ and column $q_j$, and put
\[
 U_j(X)=A_{i_j}(X)_{*,q_j},\qquad
 V_j(X)=A_{i_j}(X)_{p_j,*}.
\]
All two-by-two minors vanish, so the following is an identity
over the univariate rational function field:
\[
 A_{i_j}(X)=\frac{U_j(X)V_j(X)}{a_j(X)}.
 \tag{P14.9}
\]
Every $U_j,V_j$ is a nonzero polynomial vector, and $a_j\ne0$.
Writing $B^{(j)}$ for the product of the full-rank factors
strictly between consecutive rank-one positions, and including
the two exterior products, gives the polynomial identity
\[
 \left(\prod_{j=1}^{s}a_j(x_{\pi(i_j)})\right) f(x)
   =\bigl(uB^{(0)}U_1\bigr)
      \left(\prod_{j=1}^{s-1}V_jB^{(j)}U_{j+1}\right)
       \bigl(V_sB^{(s)}v\bigr).
 \tag{P14.10}
\]
The displayed factors on the right are scalar polynomials.
Since $f\ne0$ and the coefficient ring is a domain, each one
is nonzero. Each is a bridge of the form (P14.6), including
the case of an empty interior.

The variable at a rank-one reset appears in two adjacent
bridges, through $U_j$ and $V_j$. This causes no independence
assumption: the bridge lemma proves that each factor remains
a nonzero polynomial under the \emph{same} substitution $G$.
Their product is therefore nonzero in
$F[t,\mathbf y,\mathbf z]$. Moreover every pivot denominator
$a_j(G_{\pi(i_j)})$ is nonzero because setting all amplitudes
to zero gives $a_j(c_{\pi(i_j)}t)\ne0$. Substituting into the
polynomial identity (P14.10) now implies $f(G)\ne0$ directly.
There is no evaluation at a pole and no need to choose pivot
entries algorithmically for the black-box sampler.

If there are no rank-one factors, (P14.1) itself is one bridge
with constant endpoints and generically invertible interior.
The local lemma applies. When $w=1$ every nonzero factor is
already a scalar univariate polynomial; their product remains
nonzero on $x_i=c_it$, also covering the overlap of the two
rank alternatives in this case.

\paragraph{3. Coefficient rows, finite sampling, and rational form.}
Every coefficient row of $R(G)$ belongs to $V$, because the
substitution has coefficients in $F$. If their span were a
proper subspace of $V$, a constant column $v\in F^w$ would
annihilate all these rows while not annihilating $V$. Then
$R(x)v\ne0$ but $R(G)v=0$, contradicting the scalar result.
This proves (P14.3). Over an extension containing the nodes,
the same assertion is $V\otimes_F E$ equality over $E$.

For finite sampling, $f(G)$ has degree at most $D$ in $t$,
degree at most $(n-1)D$ in each selector $y_j$, and total degree
at most $D$ in all amplitudes $z_j$ together. Use a grid of
$D+1$ distinct center values, a grid of $(n-1)D+1$ values for
each selector, and the standard total-degree interpolation
simplex of size $\binom{D+K}{K}$ for the amplitudes. One
description of the last set is
$(b_{e_1},\ldots,b_{e_K})$ with $e_i\geq0$ and
$\sum e_i\leq D$, where $b_0,\ldots,b_D$ are distinct.
The triangular Newton basis proves its unisolvence in any
characteristic. This proves (P14.4).

For the rational form, put $A(T)=\prod_i(T-a_i)$ and
$c_i=1/A'(a_i)$. Then
\[
 \frac{\lambda_i(y)}{c_i}=\frac{A(y)}{y-a_i}.
 \tag{P14.11}
\]
Apply (P14.2) to $f(X_1/c_1,\ldots,X_n/c_n)$, whose factor
ranks are unchanged. Its argument in the original coordinates
is $t+\sum_j b_j/(y_j-a_i)$ after the dominant parameter
change $b_j=z_jA(y_j)$. These values are evaluations of a
rational function with numerator and denominator degrees at
most $K$. Conversely a generic rational function of these
degrees has a separable denominator and such a partial-fraction
expansion over an algebraic closure. Hence the two maps have
the same image closure after the stated rescaling. This
proves (P14.5); scalar contractions give its row-span version
in exactly the same way. Multiplication by a common nonzero
denominator does not alter the constant linear annihilators
of a vector of rational functions.

Finally a nonzero width-two program cannot contain an
identically zero matrix factor, and every nonzero two-by-two
factor has generic rank one or two. Thus $K=1$ applies to
all width-two programs. For $n\geq2$, the polynomial
$x_1/c_1-x_2/c_2$ has a width-two representation by the two
unipotent factors $I+(x_1/c_1)E_{12}$ and
$I-(x_2/c_2)E_{12}$, with endpoints $e_1^{\mathsf T},e_2$.
It vanishes on the common center, establishing the asserted
failure for $K=0$. $\square$

\paragraph{Scope and provenance.}
The support argument for invertible interiors and the endpoint
absorption are written out above; they are the elementary
mechanisms used in I-17 and I-30. The deduction here uses
rank-one factorization and simultaneous product nonvanishing
to replace a full matrix-image requirement by scalar bridge
conditions. It improves I-29's $r=1$ pair budget from $3w-3$
to $w-1$ under the stated factor-rank hypothesis. For a
nonzero scalar program that hypothesis also forces every
interval rank into $\{1,w\}$: an interval containing a rank-one
factor has rank at most one and cannot be identically zero,
because that would make the whole scalar output zero.

The rank-one splitting is the dimension-independent form of
the established width-two reduction in Agrawal, Gurjar,
Korwar, and Saxena (2015), author-hosted version,
\S5.4, Lemma~34, printed pp.~24--25. Their Lemma~35,
printed p.~25, gives product closure by interpolation, and
Theorem~37, printed p.~26, proves polynomial-time black-box
PIT for width two. Those complete local arguments were read.
Here endpoint absorption and coefficient-row growth give the
stated common centered generator directly, so its product
closure requires no additional interpolation of a finite set.

The rational-function interpretation uses the primary
definitions and interpolation relationship in Hu, van
Melkebeek, and Morgan (2024), Definition~1.2 and
Proposition~2.2, printed pp.~5 and 16. Their Theorem~1.9,
printed p.~10, has an additional low-degree-monomial
hypothesis; it is not being used to justify the arbitrary-degree
statement here. No literature-priority claim is made for
this particular generator or the bridge formulation.
For growing $w$, (P14.4) is $(nD+1)^{O(w)}$, not polynomial
in $w$. Arbitrary intermediate factor ranks remain outside
this proof. No purity, GRH, or approximate energy equivalence
is required.
